\cotefile{157-2}
\scanpage{18}
des foncteurs \inserted{$\unsure{c'|\varepsilon'}\colon\unsure{\mathscr{L}|\mathcal{L}}_{\unsure{m'n'|m'+n'|m',n'}}\unsure{\mathord{-}|\to}X^{\unsure{\checkmark|v}}$} dont le composé avec $\alpha'$ est le foncteur
\[
\beta\colon\varepsilon\longrightarrow X\qquad
\begin{aligned}a_0&\unsure{\mathrel{\mapstochar\relbar}|\longmapsto} x\\ b_0&\longmapsto z\end{aligned}
\tag*{$(*)$}
\]
et de $=$ la catégorie-source $\unsure{A/|A}$ \deleted{\unsure{est formée de la}} \inserted{s'interprète comme la}
sous-catégorie pleine de $\underline{\mathrm{Hom}}(\mathscr{L}_{\unsure{mn|m+n}},X)$ formée
des $c\colon\mathscr{L}_{mn}\mathord{-}X$ dont la restriction à
$\{a,b,c\}$ est le foncteur
\[
a\mathrel{\mapstochar\relbar}x,\ b\mathrel{\mapstochar\relbar}y,\ c\mathrel{\mapstochar\relbar}z\quad .
\]
Elle est contenue dans la sous-catégorie pleine
$A'$ des foncteurs dont la restriction à $\{a,b\}$ soit
$a\mathrel{\mapstochar\relbar}x,\ b\mapsto z$, \unsure{i.e.} dont le composé avec $\alpha$ soit
le foncteur $\beta\colon\varepsilon\mathord{-}X$ déjà envisagé ci-dessus.
Ainsi les foncteurs $f,g\colon A\rightrightarrows B$ se factorisent
par $A'$
\diagram{\begin{tikzcd}[ampersand replacement=\&,column sep=large]
 A \arrow[r,hook,no head] \arrow[drr,shift left=1,"f"] \arrow[drr,shift right=1,"g"'] \& A' \arrow[dr,shift left=1,"f'"] \arrow[dr,shift right=1,"g'"'] \& {}\\
 {} \& {} \& B
\end{tikzcd}},
foncteurs déduits de $\varphi,\psi$ ci-dessus, donnant
lieu \unsure{aux} triangles commutatifs \unsure{(3)}, et \deleted{\unsure{au}} foncteur
\deleted{\unsure{fonct.}} $\beta\colon\varepsilon\mathord{-}X$, de la façon dite. Pour
prouver que $f$ et $g$ sont homotopes, il suffit
de le prouver pour $f',g'$. Et pour ceci, il
\deleted{\unsure{revient au même}} \inserted{suffit} de montrer que $\varphi$ et $\psi$
sont homotopes en \unsure{tant} qu'objets de la
catégorie $\underline{\mathrm{Hom}}(\mathscr{L}_{m+n},\unsure{\mathscr{L}}_{\unsure{m'n'|m'+n'|m',n'}}\unsure{,};\alpha)$ formée
des foncteurs dont la restriction à $\varepsilon$ est $\alpha$.
% archive sheet number: 17
